\documentclass[../main.tex]{subfiles}

\begin{document}

\begin{abstract}

Abrupt permafrost thaw disproportionately accelerates carbon release and damages infrastructure, yet existing spatial products describe thaw occurrence, landform, or stage rather than thaw mode. We train a gradient-boosted classifier on 19,288 expert-labeled observations, using 70 environmental variables to distinguish thaw mode across Alaska. Because abrupt observations dominate the training data, we evaluate skill on the minority non-abrupt class and test generalization with spatial cross-validation. We convert the output into an index of how closely local conditions resemble abrupt or non-abrupt thaw sites, independent of each mode's frequency in training, and bound predictions to terrain that is quantitatively similar to training points. The classifier achieves an area under the precision-recall curve of $0.852\pm 0.011$, roughly fifteen times the $0.057$ chance floor. Within the area-of-applicability, covering about $90\%$ of Alaska's permafrost domain, $24.7\%$ of the landscape has features more consistent with abrupt than non-abrupt thaw, or 264,000 km$^2$. This terrain concentrates in the northern and western tundra, including roughly two-thirds of the Brooks Range and Seward Peninsula. Grouped Shapley (SHAP) values identify terrain characteristics and temperature regime as leading indicators of thaw mode, and the model links both flat, low-lying wetlands and steep uplands to abrupt thaw, consistent with distinct thaw pathways. Previously mapped thermokarst occurrence explains only $1-7\%$ of the variance in the index, showing that thaw mode is a largely distinct, complementary dimension of permafrost change. The resulting surface maps the landscape's relative susceptibility to abrupt thaw and should guide field campaigns, permafrost-carbon models, and regional hazard assessments.

   \begin{plainlanguagesummary}
   Permafrost is ground that remains below $0^\circ$C for two consecutive years or more, and it does not thaw uniformly everywhere. Abrupt thaw, marked by extreme and rapid changes to ecosystems and the land surface, exposes deep carbon stores and threatens human infrastructure and communities. We train a machine learning model on 19,288 observations of permafrost thaw, labeled by experts, using 70 different environmental factors to map where conditions across Alaska resemble those linked to abrupt or non-abrupt thaw. We test the model against points from regions it did not see in training, and limit its predictions to places where the landscape is similar to our training points. Within this area, about one quarter, or 264,000 km$^2$, has conditions more consistent with abrupt than non-abrupt thaw. These areas concentrate in the northern and western tundra, including much of the Brooks Range and Seward Peninsula. Topography and temperature are the most important indicators separating the two thaw modes. The map does not show where thaw is happening now, nor predict where it will occur in the future. Instead, it shows where Alaska's landscapes may be more susceptible to abrupt than non-abrupt thaw, helping guide field observations, climate modeling, and studies of regional hazards.
   \end{plainlanguagesummary}

\end{abstract}

\end{document}
